\section{Answering the Research Question}

The model is solved under both tax regimes and calibrated to a joint-taxation baseline.
The counterfactual switches the regime to individual taxation on a revenue-neutral basis.
I compare the two steady states along three dimensions: (i) the gender gap in hours worked at each age, (ii) the lifetime human capital trajectories of both spouses, and (iii) household welfare measured by the consumption-equivalent variation.

\noindent The central prediction is that individual taxation raises the secondary earner's labor supply directly through the lower effective marginal rate, and indirectly through the human capital channel: more hours today raise future wages, which further increase labor supply in subsequent periods.
The dynamic amplification separates this framework from a standard static tax analysis and is the key contribution of the proposed model.

\noindent As an additional exercise, the transition path from joint to individual taxation can be simulated to quantify adjustment costs and the speed at which the gender gap closes following a reform.
Because the human capital channel operates slowly, the long-run effect of the reform is expected to exceed the short-run impact, and the transition simulation makes this distinction concrete.

\noindent To assess heterogeneity, I solve the model separately for households at different points of the wage distribution.
Couples where the two spouses have similar initial wages face a smaller joint-taxation distortion to begin with, so the reform is expected to have a larger proportional effect on households with high within-couple wage inequality --- precisely the group where the secondary earner penalty is most binding.

\noindent Finally, the welfare analysis decomposes the consumption-equivalent variation into a within-household redistribution component and an efficiency component.
Even holding total household income fixed, a shift to individual taxation changes the relative prices faced by each spouse, potentially increasing total household surplus by reducing the over-specialization that joint taxation induces.
